\documentclass{article}
\usepackage{graphicx}
\usepackage{amsmath}
\usepackage{amssymb}
\usepackage{amsfonts}
\usepackage{hyperref}
\usepackage{exercises}
\usepackage{indentfirst}
\usepackage{circledsteps}
\usepackage{amsthm}
\usepackage{csquotes}
\usepackage{yhmath}
\usepackage{comment}
\usepackage{placeins}
\usepackage{geometry}
\usepackage{hyperref}
\usepackage{tikz}
\usepackage{tabto}
\newcommand*{\oldneg}{\mathord{\sim}}
\newcommand{\Z}{\mathbb{Z}}
\newcommand{\N}{\mathbb{N}}
\let\oldemptyset\emptyset

\theoremstyle{definition}
\newtheorem{definition}{Definition}
\newtheorem{theorem}{Theorem}[section]
\newtheorem{corollary}{Corollary}[theorem]
\newtheorem{lemma}[theorem]{Lemma}

\setlength{\parindent}{.5in}
\geometry{margin=0.5in}
\setlength{\parskip}{12pt}
\title{Chapter 16 Notes}
\begin{document}
\date{}
\maketitle

\section*{Section 1}

\subsection*{Theorem 2.7}
Here is a proof of the classical part of this theorem. Consider some model modal 
$\mathcal{M}$ for which the hypotheses of the problem hold

Atomic case: In this case $X$ is atomic. Thus it must be that $A=X$ and $Y=B$, meaning, because 
we have assumed $A\equiv B$ is valid in $\mathcal{M}$, $X\equiv Y$ (which is the same formula) is also valid in $\mathcal{M}$.

Negation case: In this case $X$ is of the form $\oldneg\phi$. Then either $A=X$ or 
$A$ is a subformula of $\phi$. In the first case, obviously it follows that $X\equiv Y$ is valid in 
$\mathcal{M}$. In the second case, $Y$ must be of the form $\oldneg\psi$ (where $\phi\equiv\psi$ is valid over 
$\mathcal{M}$ by assumption). Then $\phi\supset\psi$ is valid over $\mathcal{M}$ (by assumption), meaning 
$\oldneg\psi\supset\oldneg\phi$ is valid over $\mathcal{M}$. Similarly, 
$\oldneg\phi\supset\oldneg\psi$ is valid over $\mathcal{M}$ (because $\psi\supset\phi$ is valid over $\mathcal{M}$).
Thus $\oldneg\phi\equiv\oldneg\psi$, that is, $X\equiv Y$, is valid over $\mathcal{M}$.

Conjunction case: In this case $X$ is of the form $\phi\wedge\psi$. Then either $A=X$ or 
$A$ is a subformula of either $\phi$ or $\psi$. In the first case, obviously it follwos that 
$X\equiv Y$ is valid in $\mathcal{M}$. In the second case, $Y$ must be of the form 
$\theta\wedge\gamma$, where both $\theta\equiv\phi$ and $\gamma\equiv\psi$ are true in $\mathcal{M}$.
Then $(\phi\wedge\psi)\equiv(\theta\wedge\gamma)$, that is, $X\equiv Y$, being valid in $\mathcal{M}$ follows 
easily.

Existential case: In this case $X$ is of the form $(\exists x)\phi(x)$. Then either $A=X$ or 
$A$ is a subformula of $\phi(x)$. Again, the first case is obvious. In the second case, we show that 
(1) $(\exists x)\phi(x)\supset (\exists x)\psi(x)$ is valid over $\mathcal{M}$. So, for any world $p$ of the 
model, suppose that $(\exists x)\phi(x)$ is true at $p$. Then $\phi(a)$ is true at $p$ for some 
$a\in\mathcal{D}$. It is easy to see that this means $\psi(a)$ msut be true at $p$ (because $A\equiv B$ is assumed to be 
valid over $\mathcal{M}$). Thus $(\exists x)\psi(x)$ is true at $p$. Hence, because (1) is true at all possible worlds of 
$\mathcal{M}$, it follows that it is valid over $\mathcal{M}$. Similar reasoning shows that the converse of 
(1) is valid over $\mathcal{M}$. Hence $(\exists x)\phi(x)\equiv(\exists x)\psi(x)$, that is, 
$X\equiv Y$, is valid over $\mathcal{M}$. 

\end{document}
